%% 第二章

\chapter{相关研究与理论基础}
\chaptermark{相关研究与理论基础}

\section{无人机具身智能与低空多模态评测}

无人机具身智能研究关注智能体如何在三维空间中感知环境、理解任务并执行动作。与地面机器人相比，无人机视野范围大、运动自由度高，但是视角变化剧烈、目标尺度变化复杂、时间连续性强\cite{shao2025uavvln,wang2025embodiedintelligencesurvey,sun2025embodiedintelligentautonomousuav}。近年来，低空无人机任务逐渐从传统检测、跟踪和遥感理解扩展到视觉语言导航、城市问答、目标搜索和高层任务规划等方向。

AerialVLN 将视觉语言导航引入无人机场景，强调无人机在连续飞行中依据语言指令完成导航任务\cite{liu2023aerialvln}。LongFly、AeroDuo、OpenFly、SkyVLN 以及 See, Point, Fly 等工作进一步关注长时程空中导航、多无人机场景协同、统一导航平台、控制耦合导航和免训练导航框架，为多视角低空任务提供了新的研究基础\cite{jiang2025longfly,wu2025aeroduo,OpenFly,li2025skyvln,pmlr-v305-hu25e}。CityNav 和 CityNavAgent 面向真实城市空间构建大规模空中导航任务，强调层次化语义规划、全局记忆和长程路径理解\cite{lee2025citynav,zhang2025citynavagent}。进一步地，CityEQA、HUGE-Bench、LogisticsVLN、UrbanVideo-Bench、CityAVOS等工作，将具身问答、视觉语言动作任务、低空物流导航、城市视频具身评测和目标搜索代理纳入无人机多模态评测范围\cite{zhao2025cityeqa,guo2026huge,zhang2025logisticsvln,zhao2025urbanvideobench,ji2026towards}。

上述工作推动了无人机多模态任务的发展，但多数评测仍更强调环境理解、任务完成或导航结果。例如，模型往往需要判断目标在哪里、该如何到达目标，或者视频中发生了什么事件。但是智能体在真实无人机任务中还要知道自身的状态，比如相对地标的位置、当前朝向、未来视角、飞行行为等。如果评测只考察外部环境而忽略无人机自身状态，就难以全面刻画具身推理能力。

本文的双认知框架就是针对这一不足提出的。UAV-DualCog 把无人机自身的状态和外部环境作为顶层认知对象，在图像和视频两种媒介中同时设置任务，使得评测可以覆盖“我在哪里、我会看到什么、我正在怎样飞行”以及“目标在哪里、我应该怎样接近目标、目标何时可见”等两方面的问题。

\section{多模态大模型综合评测与空间推理}

近年来，GPT-5.5、Claude Opus 4.7、Qwen3.7 Max、MiMo-2.5-Pro等多模态大模型不断推动视觉语言模型在感知、理解和推理上的能力边界\cite{hurst2024gpt4o,wang2024qwen2vl,wang2025internvl3_5}。在此基础上，多模态大模型评测已经形成较丰富的任务体系。MME 从感知和认知两个方面评估多模态模型基础能力\cite{fu2023mme}；MMBench 采用选择题形式评估模型是否具备较全面的多模态理解能力\cite{liu2024mmbench}；SEED-Bench、MM-Vet、MM-Vet v2、MLLM-Bench 和 MMMU 等工作进一步考察复杂视觉问答、组合推理、跨学科理解和综合能力\cite{li2024seed,yu2023mmvet,yu2024mmvet2,ge2025mllm,yue2024mmmu}。这些基准推动了多模态大模型从感知型任务向推理型任务发展。

在空间推理方面，Thinking in Space、HiSpatial、Open3D-VQA 和 CityRefer 等工作开始关注模型对空间关系、三维结构、开放空间定位和跨视角关系的理解\cite{yang2025thinking,liang2026hispatial,zhang2025open3d,miyanishi2023cityrefer}；而面向空中智能体的统一框架也进一步强调空间、时间与具身推理的耦合\cite{xu2025aerial,yang2026vision}。这些研究说明，模型不应只识别目标“是什么”，还应进一步说明目标“在哪里”、该判断基于怎样的几何关系以及模型是否具备跨视角空间记忆。

但是目前大多数综合评测都是用静态图像或者短视频作为输入，视觉内容和智能体的主动运动关系比较弱。在这样的被动观察中，模型即使得到高分，也不能直接说明它在无人机主动飞行、视角快速变化、目标可见性不断变化的情况下仍然可以保持稳定。更重要的是，很多评测只看答案是否正确，而不会要求模型给出可检验的空间或者时间证据。UAV-DualCog 因此将证据纳入正式评测，图像任务除了要输出选项外还要输出目标框，视频任务除了要输出行为或者次数外还要输出对应的时间区间。

\section{能力优化训练方法}

多模态大模型能力优化一般用后训练的方法来完成。监督微调是使用高质量样本使模型学习任务输入到结构化输出的映射关系，是目前使用最广、最稳定的微调方式。LoRA 和 QLoRA 等参数高效微调方法降低了微调资源需求，使中小规模多模态模型能够在有限算力条件下完成任务适配\cite{hu2021lora,dettmers2023qlora,han2024peft}。

对多模态结构化证据任务来说，监督微调的意义就是使模型学会稳定的输出答案和证据的形式。图像任务的多项选择答案和目标框可以转换成文本监督目标，视频任务的语义结果和时间区间也可以用结构化的JSON表达出来。这样模型就先建立基本的任务格式以及语义与证据之间的关系。

但是已有研究以及实验现象也表明，仅仅依靠监督微调并不能很好地解决细粒度证据一致性问题。模型可以输出格式正确的答案，但是仍然不能依靠正确的目标框或者时间区间进行推理。近些年来，基于可验证奖励的强化后训练方法受到重视。GRPO 等组相对策略优化方法能够在不引入复杂 critic 的情况下，利用格式奖励、答案奖励和证据质量奖励对模型进行小规模校正\cite{shao2024deepseekmath,wen2025verifiable}。

本文能力优化部分是在上述后训练方法的基础上建立起来的，先用 UAV-DualCog-Train 对基座模型做多模态监督微调，使其学会稳定的答案加证据的结构化输出，然后用联合 GRPO 同时加入格式、语义、空间证据奖励来对语义-证据一致性进行定向优化。
\section{本章小结}
本章从无人机具身智能、低空多模态评测、多模态大模型综合评测、空间推理、能力优化训练方法等几个方面对相关研究进行了梳理。现有的无人机研究已经涵盖了视觉语言导航、城市空间问答、目标搜索以及低空视频理解等一系列任务，为创建无人机多模态测试基准打下了重要的基础；目前有关多模态评价及空间推理的研究，又促使模型由感知型任务向关系理解以及复杂推理任务发展。但是大多数工作仍然以外部环境理解、导航结果或者通用视觉问答为主，缺少将无人机自身状态建模和外部环境理解并列起来的统一评测框架，也较少把目标框和时间区间等显式证据作为正式评测对象。

围绕能力优化，已经存在监督微调、参数高效微调以及基于可验证奖励的强化后训练等方法为本文提供技术基础。但是无人机双认知任务要求模型同时保证语义答案正确、空间或者时间证据可靠、跨媒介表现稳定，所以不能直接使用通用后训练的目标。本文后面几章在此基础上展开，第三章建立面向自我状态认知、环境情况认知的无人机双认知测试基准并分析模型能力结构，第四章根据评测结果来设计证据感知的训练集以及能力优化的方法。
